\documentclass[12pt]{article}
% Préambule commun à tous les documents extraits de « La Revue CAPES »
\usepackage[T1]{fontenc}
\usepackage[utf8]{inputenc}
\usepackage{lmodern}
\usepackage{amsmath,amssymb,amsthm,amsfonts,mathrsfs}
\usepackage{setspace}
\usepackage{listings}
\usepackage{pgfplots}
\pgfplotsset{compat=1.18}
\usepackage{longtable}
\usepackage{adjustbox}
\usepackage{lettrine}
\usepackage{tkz-tab}
\usepackage{changepage}
\usepackage{pdflscape}
\usepackage{graphicx}
\usepackage{tikz}
\usepackage{xcolor}
\usepackage[a4paper,top=2cm,bottom=2.5cm,left=2.5cm,right=2cm]{geometry}
\usepackage{fancyhdr}
\usepackage{draftwatermark}
\SetWatermarkText{La RevueCapes}
\SetWatermarkLightness{0.9}
\SetWatermarkScale{2}
\usepackage{hyperref}
\hypersetup{colorlinks=true,linkcolor=blue!50!black,urlcolor=blue!50!black}

\definecolor{revue}{RGB}{20,60,120}

\pagestyle{fancy}
\fancyhf{}
\renewcommand{\headrulewidth}{0.4pt}
\setlength{\headheight}{15pt}

% \entete{Matière}{Titre}{Type de document}
\newcommand{\entete}[3]{%
  \fancyhead[L]{\small\textsc{La Revue CAPES} -- #1}%
  \fancyhead[R]{\small #2}%
  \fancyfoot[C]{\thepage}%
  \thispagestyle{plain}%
  \begin{center}
    {\color{revue}\rule{\textwidth}{1.2pt}}\\[4pt]
    {\small\textsc{Concours directs CAP-CEG / CAPES -- Burkina Faso}}\\[6pt]
    {\LARGE\bfseries\color{revue} #2}\\[6pt]
    {\large\itshape #1 -- #3}\\[2pt]
    {\color{revue}\rule{\textwidth}{1.2pt}}
  \end{center}
  \vspace{0.5cm}%
}

% Encadré pour les remarques ajoutées dans les corrigés
\newcommand{\remarque}[1]{\par\medskip\noindent\fcolorbox{revue}{revue!6}{\parbox{\dimexpr\linewidth-2\fboxsep-2\fboxrule}{\textbf{\color{revue}Remarque.} #1}}\par\medskip}


\begin{document}
\entete{Culture générale}{Méthodologie de la dissertation de culture générale}{Cours}
\begin{spacing}{1.5}

 Une \textbf{dissertation} est la formalisation d'un raisonnement fondé sur la définition et la position d'un problème (Introduction) et aboutissant à la proposition d'une solution (B et/ou C de la troisième partie).

 Une dissertation doit mener un raisonnement progressif, construit, qui réponde 
à une problématique dégagée d'un sujet posé, presque toujours, sous la forme d'une 
citation ou d'une question. Elle comporte trois parties principales : une introduction; un plan progressif ( ou développement) généralement en trois parties, chacune composée de plusieurs sous-parties, des exemples bien développés et judicieux pour étayer le raisonnement; et une conclusion.
 
 
 \textbf{  L'évaluation privilégiera :}
 
\begin{itemize}
 \item la pertinence de l'analyse de la question dégagée à partir du sujet ;
 \item l'intérêt et la qualité de la réponse apportée à cette question initiale ;
\item la clarté de la démonstration, voire sa finesse, et sa progression ;
 \item la solidité de la culture littéraire du candidat(e) ;
 \item la justesse, la variété et la qualité de développement des exemples proposés pour illustrer la démonstration ;
 \item la qualité de l'expression (correction grammaticale et lexicale, orthographe).
\end{itemize}

\section*{B. Les différentes parties d'une dissertation}

Une dissertation comporte trois parties :

\quad
\textbf{ 1. L'introduction}

 Elle se fait en trois temps :

\quad

\textbf{a. Préambule ou présentation de l'œuvre}


Introduction du thème littéraire général dans lequel s'inscrit le sujet (la lecture, le roman, l'écriture, la poésie,etc.). Il s'agit d'entrer dans le devoir en justifiant le sujet/présentant le thème, le domaine auquel il se rattache.\\
Si vous trouvez rapidement une phrase d'ordre général (ou phrase dite d' $\ll$ accroche $\gg$) qui serve d"entrée en matière, c'est mieux. Mais ne perdez pas de temps à chercher une 
accroche séduisante, cela a peu d'importance. La phrase d'accroche peut concerner l'objet d'étude, le mouvement littéraire ou le thème.

Les phrases suivantes doivent présenter l'œuvre (ou les œuvres) avec sa 
date de parution, son auteur et le contexte littéraire
 
\textbf{ Il faut éviter de commencer par :}

$\circ$ Des platitudes trop générales, telles que $\ll$ depuis que l'homme est homme$\gg$, $\ll$ de tous temps $\gg$, $\ll$Depuis la nuit des temps$,\gg$ $\cdots$ ;

 $\circ$  Des remarques pleines de jargon qui veulent impressionner le correcteur ;

$\circ$ Des élucubrations inutiles sur la vie de l'auteur de la citation.\\ 



 
 
\textbf{\underline{ Attention :} Si vous commencez par une autre citation, il faut être prudent : elle 
doit être très $\ll$ accrocheuse $\gg$ mais ne doit pas créer de confusion avec celle qui est l'objet réel de la dissertation.}

\textbf{ b. Annonce du sujet}

 Présenter le sujet précis du devoir à l'intérieur de ce thème plus vaste (telle 
fonction de la lecture, telle définition du roman, telle conception de l'écriture 
ou de la poésie).

\begin{itemize}
\item Il s'agit de présenter le sujet en précisant le nom de l'auteur (ne dites de lui que ce qui est utile pour cerner et comprendre la citation) et époque : le recopier intégralement s'il est court ; s'il est assez long, le découper en segments que vous devez analyser tout en faisant apparaître les termes-clefs de la citation entre guillemets. Définir les termes-clefs afin d'amener la problématique que vous devez formuler (qui n'est pas la répétition du sujet sous une forme 
interrogative).
\item  Éviter les séries de questions directes.
\end{itemize}

\quad
\textbf{ c. Élaborer un problématique}

$-$ Lisez plusieurs fois la reformulation du sujet rédigée à partir de vos définitions. Au brouillon, écrivez toutes les idées qui vous viennent à l'esprit sur le sujet (exemples, autres, événements, ...).


$-$ Reformulez la question posée avec vos propres mots et montrez sa richesse et ses sous-entendus pour trouver la problématique qui est souvent un paradoxe qu'il faut dépasser.

$-$ La problématique est le fil directeur du devoir : votre plan doit répondre à la question qu'elle 
pose.

Comment les termes-clefs se relient-ils, s'articulent-ils entre eux ? Jouent-ils les uns par rapport aux autres de manière paradoxale, polémique ou problématique ? Où se 
situe la tension, la contradiction, le paradoxe, le problème ?
 
 Bien découper les propositions, les articulations, le jeu entre les termes dans des 
citations longues : penser à bien prendre en compte tous les termes et les liens logiques.


 Se poser diverses questions-clefs :
 
 \begin{itemize}
 \item Quels sont les enjeux exacts de ce sujet ?
 \item Quel est l'intérêt de la citation ?
 \item Quelles sont ses limites ?
 \item Qui parle et de quel point de vue ? Est-ce un auteur ? Un personnage littéraire ? Un critique ?
 
 Il s'agit de bien distinguer le thème (ce dont l'auteur parle) de la thèse (ce qu'il en dit).
 La problématique est-elle simple (question unique) ou complexe (entrelacs de 
plusieurs questions) ?
 
 Y a-t-il un contexte précis (historique, culturel, littéraire, idéologique $...$) à prendre en compte dans la citation ? Le nom de l'auteur, l'œuvre d'où est tirée la citation et sa date 
sont généralement donnés : en tirer parti (quel sens prend cette position par rapport à 
sa propre œuvre ?) sans non plus faire une dissertation d'histoire, d'histoire littéraire ni proposer une réflexion portant exclusivement sur l'auteur et son œuvre.
 
 Y a-t-il une légère distorsion entre la citation du sujet et la question qui l'accompagne, concernant le sujet lui-même ou le champ des exemples ?
 
 Y a-t-il un ton particulier (ironique, caricatural, polémique) qui sera à prendre en compte ?
\end{itemize}

\quad

\textbf{d.  Élaborer un plan}

Le plan d'une dissertation peut prendre diverses formes. L'important est qu'il répondre bien à votre problématique pour que vous évitez le hors-sujet.

\begin{itemize}
\item Utilisez votre brouillon initial sur lequel vous avez noté vos idées.
\item Classez ensuite ces idées par thématique ou argument.
\item Normalement, vous pourrez arrivez à deux ou trois idées principales, divisées en deux ou trois sous-parties qui seront illustrées par des exemples concrets.
\item N'oubliez pas de rédiger une transition entre chaque grande partie (conclusion de la partie actuelle et introduction de la partie suivante).
\end{itemize}

\vspace*{1.5cm}

\textbf{Les différents types de plan}

En dissertation, un plan bien structuré est essentiel pour la clarté et la cohérence de l'argumentation. Les types de plans peuvent varier en fonction du sujet, du style et de l'objectif de la dissertation.
 
On distingue \textbf{ quatre grands} types de plans : 
 
 $\rhd$ le \textbf{plan dialectique :} Ce type de plan est basé sur une opposition entre deux idées ou thèses et se développe en trois étapes : thèse (oui) – antithèse (non) – synthèse ou dépassement (mais). Ce plan peut correspondre aux sujets du type « Pensez-vous que... ? » / « Dans quelle mesure ? » ; 

\textbf{Exemple : La liberté d'expression est-elle absolue ?}

\begin{itemize}
\item \textbf{Thèse :} La liberté d'expression doit être absolue.
\item \textbf{ Antithèse :}La liberté d'expression doit être limitée.
\item \textbf{ Synthèse :}La liberté d'expression doit être encadrée par des lois pour protéger les individus et l'ordre public.\end{itemize}




$\rhd$ le \textbf{plan analytique :} Dans un plan analytique, le but est de diviser le sujet en plusieurs sous-parties logiques et d'examiner chaque aspect du problème. Il est souvent utilisé pour les sujets qui exigent une analyse détaillée de différents éléments ou phénomènes. Il se divise généralement en deux ou trois grandes parties, chaque partie étant une analyse détaillée d'un aspect particulier du sujet.
Il met en lumière les différents facteurs ou causes d'un problème, sans nécessairement chercher une opposition.

\textbf{Exemple :}  L'impact des réseaux sociaux sur la société

\textbf{I. L'impact positif des réseaux sociaux}

$-$ Facilitation de la communication et de l'information.

$-$ Renforcement des liens sociaux et communautaires.

\textbf{II. Les effets négatifs des réseaux sociaux
}

$-$ Isolement social et addiction.

$-$ Diffusion de fausses informations et manipulation.

\textbf{III. Les enjeux et régulations des réseaux sociaux}

$-$ Besoin de régulation des contenus.

$-$ Protection de la vie privée et des données personnelles.

\quad




$\rhd$ \textbf{ Le Plan Comparatif}

Ce plan consiste à comparer deux ou plusieurs éléments, idées ou phénomènes afin de dégager leurs différences ou similarités. Il est souvent utilisé dans des sujets où l'on doit confronter des théories, des événements, ou des situations.

\textbf{Exemple : }Comparer le système éducatif du Burkina Faso et celui d'un autre pays africain

\textbf{I. Le système éducatif au Burkina Faso}

$-$ Structure et caractéristiques du système.

\textbf{II. Le système éducatif d'un autre pays africain}

$-$ Comparaison des approches, structures et méthodes.

\textbf{III. Comparaison des résultats obtenus et des défis}

$-$ Analyse des points forts et faibles des deux systèmes.



$\rhd$ le \textbf{plan thématique :} il propose une réponse organisée autour des différents aspects d'un thème ou d'une notion. Vous pouvez poser ici les questions Quoi ? Comment ? Pourquoi ? qui vous aideront 
à circuler dans l'œuvre.

 \vspace*{1.5cm}
 
 En plus de ces plans, ils existent d'autres types de plans comme : \textbf{le plan problématique, le plan chronologie, le plan causes/conséquences, etc. }

\underline{\textbf{Élaborer un plan détaillé.}}
 
$\circ$ Noter au brouillon des citations ou des passages qui viendront illustrer les arguments. 

$\circ$ Dégager parmi toutes les remarques les grandes parties du plan qui ira du plus évident au 
plus complexe (deux ou trois). 

$\circ$ Pour chaque partie élaborer un sous plan : de deux à trois arguments ou sous-parties. Chaque 
argument devra s'appuyer sur au moins un exemple précis, auquel peuvent s'ajouter d'autres exemples plus ou moins développés. 

$\circ$ Rédiger au brouillon les transitions qui doivent à la fois conclure la partie et amener la partie suivante. 

\quad 
\textbf{\underline{ Remarque :} Qu'est-ce qu'un bon plan ?}
 
\quad Le plus difficile est bien évidemment de trouver un plan. Il ne faut pas viser la perfection qui n'existe pas. 

\quad Le bon plan est celui qui permet de développer une véritable réflexion à partir du sujet, de progresser dans une argumentation, de développer différentes remarques sans se répéter. L'autre écueil à éviter est le hors-sujet, n'hésitez pas à rappeler à chaque début les termes du sujet et à 
montrer en quoi chacune de vos parties y répond.


\quad
\quad
\textbf{2. Le développement}

Le développement d'une dissertation comporte toujours deux ou trois parties. Si vous faites une dissertation en deux parties, vous devez rédiger trois sous-parties pour chacune (deux si vous faites trois grandes parties).

Chaque partie soutient une idée centrale qui répond à la problématique, alors que chaque sous-parties s'articule autour d'un argument qui soutient et illustre l'idée directrice.

Vos arguments doivent absolument être illustrés par un exemple !

Entre chaque partie vous devez rédiger une transition qui conclut la partie précédente et annonce la partie suivante.


\textbf{3. La conclusion}

Elle comporte :
 
$\bullet$ Une réponse à la problématique grâce à un résumé du développement qui rappelle les grandes lignes de la démarche argumentative. 

$\bullet$ Une ouverture vers des éléments de réflexion en lien avec le sujet de dissertation. Il faut éviter à tout prix de reposer une nouvelle question.

\end{spacing}
\end{document}
